\documentclass[12pt]{article}
\usepackage{graphicx}
\usepackage[margin=2cm, a4paper]{geometry}
\usepackage{setspace}
\usepackage{pdfpages}
\usepackage{float}
\usepackage{amsmath}
\usepackage{fancyvrb}
\usepackage{amssymb}
\usepackage{minted}
\usepackage{enumitem}
\usepackage[colorlinks,linkcolor=blue]{hyperref}


\renewcommand{\contentsname}{Contents}
\renewcommand{\figurename}{Figure}
\renewcommand{\tablename}{Table}
\hypersetup{
    colorlinks=true,
    linkcolor=black,
    filecolor=magenta,      
    urlcolor=blue,
}
\newcommand{\mytitle}{Introduction to Algebra (II) HW5}
\newcommand{\myauthor}{B13902022 Richard Lai}

\usepackage{fancyhdr}
\pagestyle{fancy}
\fancyhead{}
\fancyhead[L]{\mytitle}
\fancyhead[R]{\myauthor}

\usepackage[english]{babel}
\newtheorem{theorem}{Theorem}[section]
\newtheorem{corollary}{Corollary}[theorem]
\newtheorem{lemma}[theorem]{Lemma}

\usepackage{amsthm}
\theoremstyle{definition}
\newtheorem{definition}{Definition}[section]
\theoremstyle{remark}
\newtheorem*{remark}{Remark}

\title{\mytitle}
\author{\myauthor}
\date{\today}

\begin{document}
\maketitle

\section{}
Let $\zeta_n$ be a primitive n-th root of unity and $K_n=\mathbb{Q}(\zeta_n)$. Since the minimal polynomial of $\zeta_n$ over $\mathbb{Q}$ is the cyclotomic polynomial $\Phi(x)$, which has degree $\varphi(n)$. 

$\mathbb{Q}(\zeta_n)$ contains all roots of $\Phi(x)$, so it's a splitting field, hence normal. Since $\mathbb{Q}$ is characteristic 0 field, $\mathbb{Q}(\zeta_n)/\mathbb{Q}$ is separable. Here we have $\mathbb{Q}(\zeta_n)/\mathbb{Q}$ is Galois.

$G=\mathrm{Aut}(\mathbb{Q}(\zeta_n)/\mathbb{Q})$ acts on the roots of $\Phi(x)$ transitively, and every automorphism $\sigma\in G$ is determined by its action on $\zeta_n$. Since $\sigma(\zeta_n)$ must again be a primitive n-th root of unity, we have:
$$\sigma(\zeta_n)=\zeta_n^k\text{ for some }k\in(\mathbb{Z}/n\mathbb{Z})^\times$$

This gives a well-defined injective homomorphism:
$$\pi:\mathrm{Aut}(K_n/\mathbb{Q})\hookrightarrow (\mathbb{Z}/n\mathbb{Z})^\times,\ \sigma_k\mapsto k$$
where $\sigma$ maps $\zeta_n$ to $\zeta_n^k$ for some $k\in(\mathbb{Z}/n\mathbb{Z})^\times$ (The homomorphy is easy to check. $\pi(\sigma_a\circ\sigma_b)=\pi(\sigma_{ab})=ab=\pi(\sigma_a)\pi(\sigma_b)$, $\circ$ stands for composition here)

Because both $\mathrm{Aut}(K_n/\mathbb{Q})$ and $(\mathbb{Z}/n\mathbb{Z})^\times$ has the order of $\phi(n)$, this is an isomorphism.

\section{}
Let $f(x)=x^{2^n}+x+1$ be defined over $\mathbb{F}$. Assume it's irreducible.

Let $\alpha$ be a root of $f(x)$ in the splitting field of $f$, by definition we have:
\begin{align*}
    \alpha^{2^n}+\alpha+1=0 \\
    \alpha^{2^n}=-\alpha-1=\alpha+1 \\
    (\alpha^{2^n})^{2^n}=(\alpha+1)^{2^n} \\
    \alpha^{2^{2n}}=\alpha^{2^n}+1=-\alpha=\alpha
\end{align*}
The relation $\alpha^{2^{2n}}=\alpha$ guarantees that $\alpha$ is a element of $\mathbb{F}_{2^{2n}}$.

$[\mathbb{F}_{2^{2n}}:\mathbb{F}_2]=2n$, since we assume that $f$ is irreducible, the degree of $f$ must divide $2n$:
$$2^{n}\mid 2n$$

When $n=1,2$, the relation holds. But when $n>3$, $2^n>2n$ is always true (This can be proved by induction), so $2^n \nmid 2n$ when $n>3$. In conclusion, $f$ is irreducible only when $n=1$ or $n=2$.

\end{document}
% how to display codes?
% \begin{minted}[frame=lines,framesep=2mm,baselinestretch=1.2,linenos,breaklines]{python}
% \end{minted}

% how to display images?
% \begin{figure}[H]
%     \centering
%     \includegraphics[width=0.5\linewidth]{}
%     \caption{Caption}
% \end{figure}
% test
